\section{Screening molecules against illnesses}
\label{sec:illness}

The materials track answers "what properties does this composition have?". The same protocol,
applied to molecules, lets us ask a question that sounds very different but is structurally
identical: \emph{given a molecule, which illnesses is it worth testing against?} This section
explains why that question can be answered at all, what is actually being predicted, and the four
places where an honest measurement differs sharply from a flattering one.

\subsection{What is predicted, stated precisely}

No model in this project predicts that a molecule cures a disease. Such a model cannot be trained,
because the label does not exist: there is no dataset of \emph{molecules that cure diabetes}. What
\emph{does} exist, in enormous quantity, is measured affinity between a molecule and a
\textbf{single protein} --- the classical QSAR label, and the same kind of label used by
established tools such as SwissTargetPrediction.

Illnesses enter only as \emph{context}. Public biomedical databases record, for a given illness,
which proteins the evidence associates with it and how strongly. Combining the two gives the
chain actually implemented:

\begin{center}
\fbox{\parbox{0.93\textwidth}{\centering
molecule $\longrightarrow$ \textbf{one QSAR model per protein} $\longrightarrow$ activity profile
over \TgtNModels{} proteins $\longrightarrow$ \textbf{weighted by association evidence}
$\longrightarrow$ a score for every illness}}
\end{center}

Everything the interface reports is a predicted interaction with a protein, together with the
public evidence linking that protein to an illness. It is in-silico screening for research and
teaching --- a hypothesis about what is worth measuring next, never a claim about treating anyone.

\subsection{Why one model per protein, not one per illness}

This is the single design decision the whole module rests on, and it is worth stating separately
because three useful properties follow from it.

\begin{enumerate}
  \item \textbf{Asking about a second and third illness is free.} A molecule is featurised once
    and passed through the \TgtNModels{} protein models once. Each illness is then a different
    weighted lookup over the \emph{same} activity profile (Equation~\ref{eq:illness-score})
    --- arithmetic, not another model. The
    complete ranking over all \TgtNDiseases{} illnesses is returned by a single request, which is
    exactly the interaction the application needs: \emph{and what about this illness? and this
    one?}
  \item \textbf{Adding an illness costs no training at all.} A new illness requires one database
    query for its protein list. If those proteins already have models, the illness becomes
    screenable immediately.
  \item \textbf{The intermediate quantity is the interpretable one.} "Predicted active against
    ACHE and BCHE" is a chemically checkable statement that a domain expert can accept or reject.
    "Might help Alzheimer's disease" is not.
\end{enumerate}

\subsection{Data}

The knowledge base is built from the Open Targets GraphQL API for \TgtNDiseases{} well-studied
illnesses: their most strongly associated proteins with an association score, and the
small-molecule drugs already approved for each. Of \TgtNProteins{} associated proteins,
\TgtNModels{} have enough measured activity in ChEMBL to model, giving \TgtNMols{}
molecule--protein measurements in total.

Activity labels come from pChEMBL (the negative logarithm of the measured potency):
\textbf{$\geq 6.5$ active, $< 5$ inactive, and the uncertain band between the two discarded.}
Published benchmarks of ChEMBL show there is no standard cutoff and that the resulting
active/inactive ratio swings strongly with the choice, so the cutoff is recorded in every
per-protein summary rather than left implicit. Duplicate molecule--protein pairs are collapsed to
their median value, and all SMILES pass through the same canonicalisation used elsewhere in the
molecular track. Features and splits are unchanged from the toxicity models: RDKit descriptors
plus a Morgan fingerprint, and splits grouped by Bemis--Murcko scaffold.

\subsection{Six ways to get a flattering number, and what was done instead}

Each of the following produced a result that looked good before it was corrected. They are
reported here because the corrections, not the headline scores, are the actual contribution of
this module. The last one was not a measurement problem at all but a silent data error, found
only because a metric came out implausibly high.

\paragraph{1. Assumed inactives inflate the score.} The published literature reports what
\emph{binds} far more often than what does not, so for many proteins ChEMBL supplies plenty of
actives and too few measured inactives. Where that happened, property-matched molecules from the
unlabelled ZINC pool were added as presumed inactives (matched on molecular weight and
lipophilicity), averaging \TgtDecoyFrac\% of rows. But a ZINC molecule is \emph{assumed} inactive,
not measured inactive, and it is easy to separate from a potent drug-like active --- so it
flatters the model.

The response is not to hide those rows but to declare them. Every model reports its decoy
fraction, and ROC-AUC is reported \textbf{twice}: over all held-out molecules, and over
\textbf{measured molecules only}. The difference is the honest cost of the assumption, and for
individual proteins it is severe: in Table~\ref{tab:targets} DRD2 falls from $0.989$ to $0.792$
and SLC6A4 from $0.993$ to $0.845$. Reporting only the first column would have produced a table
of near-perfect scores that mean very little.

\paragraph{2. A measured-only score is only as good as its denominator.} Correcting for presumed
inactives immediately raises a second question that is easy to skip past: how many measured
molecules is that second number computed against? For NR3C1 the answer is \textbf{fifteen}. Of
$1{,}112$ measured molecules, fifteen are inactive --- so its measured-only ROC-AUC has almost
nothing to rank against, and it duly comes out at $1.000$. That is not a perfect model; it is the
absence of a measurement wearing the costume of a perfect result.

Rather than quietly drop such targets, which would hide information a reader can weigh perfectly
well, the denominator was made part of the claim. Every measured-only score now travels with the
number of measured molecules of its rarer class --- the \emph{measured minority} column of
Table~\ref{tab:targets}. Below fifteen the score is withheld outright as noise, which happens for
\TgtNWithheld{} of the \TgtNModels{} targets --- EIF2AK4 is the extreme case, with
\textbf{no measured inactives whatsoever}, so its entire negative class is presumed. At or above
\TgtMinMinority{} the target is called \emph{well supported}, and only those
(\TgtNWellSupported{} of \TgtNModels) enter the project-wide average: \TgtAUCMeasured{} on
measured molecules against \TgtAUC{} over all of them. The \TgtNThin{} targets in between are
reported and marked with a dagger, so that DRD2's $0.792$ is not quoted as though it had been
measured on thousands of molecules when it rests on $55$.

Table~\ref{tab:targets} should be read with that column first. ATM is its most trustworthy row:
no presumed inactives at all, $532$ measured molecules of the rarer class, and the same $0.981$
whichever way it is computed. The rows that still reach $1.000$ --- MC4R and NR3C1 --- are the
ones to distrust, not the ones to celebrate. The pattern across the table is consistent: the
larger the measured minority, the smaller the gap between the two ROC-AUC columns, which is
exactly what one expects if the gap is caused by the presumed inactives rather than by the
models.

\maybeinput{targets_table}

\paragraph{3. Enrichment measured against the test set is arithmetically meaningless.} The first
version of the enrichment factor asked whether known actives rank above the \emph{other test
molecules}. With roughly 70\% actives in a test set, the top 1\% can be at most $1/0.7 = 1.43$
times as active as the base rate, so the metric is incapable of distinguishing a good model from a
mediocre one. Enrichment is now measured the way a real virtual screen works: known actives
against a background of \textbf{20{,}000 random ZINC molecules}, which is what the top of a screen
actually competes against. Values of $\TgtEFMin\times$ to $\TgtEFMax\times$ (mean
$\TgtEF\times$) now carry information --- and the low end carries it too: GBA1, the weakest model
in the set at a measured ROC-AUC of $0.843$, enriches only $\TgtEFMin\times$, which is the
pipeline correctly declining to promise much from $451$ molecules.

\paragraph{4. Scaffold leakage.} A protein's known actives are frequently analogues of one
another, published as a series around one chemical core. A random split therefore puts near-copies
on both sides of it and measures memory rather than generalisation. Splits are grouped by
Bemis--Murcko scaffold, as in the toxicity models, so an entire chemical series falls on one side.

\paragraph{5. A fixed applicability-domain cutoff is arbitrary.} The usual practice is a hard
Tanimoto threshold of $0.4$ to the nearest training molecule. Rather than adopt that number,
held-out molecules are binned by similarity to their nearest known active and accuracy is measured
within each bin, so the unreliable region is placed where accuracy actually degrades. A molecule
that falls outside the domain is reported as \textbf{unclear}, not as a confident low score ---
extrapolation is not a finding, and the interface is explicit about the difference.

\paragraph{6. The model was right; the protein was wrong.} The knowledge base names proteins by
gene symbol, and those symbols have to be resolved to ChEMBL target identifiers before any
activity can be downloaded. The first version of that resolver searched ChEMBL by name and, if no
target carried the gene symbol, fell back to the best-scoring human single protein. For one gene
out of \TgtNProteins{} that fallback fired: \textbf{FLG} (filaggrin, a structural skin protein
associated with asthma) was resolved to \textbf{CHEMBL3650, which is FGFR1} --- a receptor
tyrosine kinase, an entirely different protein. ChEMBL duly returned $6{,}592$ real activities,
the pipeline trained a perfectly good model on them, and it was reported under the name
\emph{filaggrin} with a ROC-AUC of $1.000$.

Nothing in the training or evaluation code could have caught this: every number was correct for
the data it was given. It surfaced only because a ROC-AUC of exactly $1.000$ on a protein with no
known small-molecule chemistry was implausible enough to check, after which all
\TgtNProteins{} resolutions were audited against the gene symbols ChEMBL itself records
--- one was wrong, the other \TgtNProteins{} were correct. The resolver now requires an exact
gene-symbol match and \textbf{drops a target rather than guessing}, the mis-resolved target and
its model were deleted, and asthma retains its three correctly resolved proteins (ADORA1, PDE4D
and NR3C1, the glucocorticoid receptor that inhaled corticosteroids act on).

The lesson is the one worth carrying out of this project: a silent identifier error is invisible
to every validation metric, because the metrics never see the name. An implausibly good score is
a reason to audit the inputs, not to celebrate.

\subsection{Composing protein predictions into an illness score}
\label{sec:illness-score}

An activity profile gives one probability per protein. Turning that into a score for an illness
means choosing how to combine the proteins associated with it, and that choice turned out to
matter more than any modelling decision in this module.

\paragraph{The first attempt, and the molecule that refuted it.} The obvious choice is the
evidence-weighted mean over the illness's proteins,
$\sum_t w_{dt} p_t / \sum_t w_{dt}$. It was implemented, and it fails on the first drug anyone
would test it with.

\textbf{Donepezil} is the standard treatment for Alzheimer's disease and a selective
acetylcholinesterase inhibitor. Screened against the \TgtNDiseases{} illnesses, the pipeline put
its real target ACHE at $99\%$ with a Tanimoto of $1.00$ to a known active --- the chemistry was
recognised perfectly. Alzheimer's is associated with four modelled proteins of comparable
evidence, and donepezil is correctly predicted inactive on the other three:

\begin{center}\small
\begin{tabular}{lccc}
\toprule
Protein & evidence $w$ & $P(\text{active})$ & $w \cdot p$\\
\midrule
ACHE  & 0.75 & 99\,\% & \textbf{0.74}\\
ACE   & 0.70 &  6\,\% & 0.04\\
APP   & 0.68 &  1\,\% & 0.01\\
CDK5  & 0.60 &  4\,\% & 0.02\\
\midrule
\multicolumn{3}{l}{evidence-weighted mean} & 0.25\\
\multicolumn{3}{l}{best evidence-weighted target} & \textbf{0.74}\\
\bottomrule
\end{tabular}
\end{center}

The mean gave Alzheimer's \textbf{0.25} and the verdict \emph{unlikely}, ranking it below
depression. The pipeline had identified the correct mechanism of the correct drug and then
reported it as not worth testing, because three proteins that no drug hits dragged the average
down.

The error is conceptual, not numerical. A mean asks \emph{does this molecule engage the whole
panel of proteins linked to this illness?} Almost no drug does, and the ones that come closest
are promiscuous compounds that fail for that very reason. Selectivity is how a drug is supposed
to work, and the mean penalises it hardest.

\paragraph{What the score is now.} The question worth asking is whether \emph{any} well-evidenced
protein is likely to be engaged, so the score is the best evidence-weighted target:

\begin{equation}
  S(m, d) \;=\; \max_{t \in T_d} \; w_{dt} \cdot P(\text{active} \mid m, t)
  \label{eq:illness-score}
\end{equation}

where $w_{dt}$ is the association score between illness $d$ and protein $t$, and $T_d$ is the set
of its proteins \emph{that have a model}. Under this definition donepezil scores $0.74$ on
Alzheimer's, \emph{worth testing}, attributed to ACHE. The evidence-weighted mean is still
computed and returned alongside, as \texttt{panel\_mean}: ``engages the whole panel'' is a real
question, it is simply not the one being asked.

\paragraph{Why not a noisy-OR.} The other natural aggregation is
$1 - \prod_t (1 - w_{dt} p_t)$, the probability that at least one target is engaged, which has the
appeal of rewarding a molecule that hits several proteins. It was rejected, because it assumes the
per-target predictions are \textbf{independent} --- and they are not. Every $p_t$ is computed from
the same descriptors and the same fingerprint of the same molecule, so the predictions are
strongly correlated by construction. A noisy-OR over four correlated mediocre predictions returns
a confident-looking score that no single piece of evidence supports: four proteins at $p = 0.3$
with $w = 0.7$ would yield $0.61$, enough for \emph{worth testing} on no good grounds. The maximum
makes the weaker claim, and it is the claim the data supports.

\paragraph{A missing model is not a zero.} $T_d$ ranges only over proteins that have a model. A
protein without one is \textbf{excluded rather than counted as inactive}, because a missing model
is missing information. Each illness instead reports its \emph{coverage}, the fraction of its
associated proteins that are modelled, so a low score for a poorly covered illness cannot be
mistaken for a prediction that the molecule does nothing.

\paragraph{Confidence follows the driver.} Under a mean every protein contributed, so the least
reliable applicability domain governed the whole result. Under Equation~\ref{eq:illness-score} the
claim rests on one protein, so confidence is now read from \emph{that} protein's similarity to its
nearest known active. This is why donepezil's Alzheimer's result is reported at \textbf{high}
confidence --- the driving target is an exact structural match --- where the previous scheme
reported \emph{low}, having been dragged down by the unrelated proteins it was averaging over.

The score never appears alone. It is returned with the per-protein probability, the association
weight, the nearest known active and its similarity, the driving protein, the coverage and a
confidence flag. Toxicity flags from the Tox21 models then downgrade an otherwise promising
result, and the verdict is reported in three bands: \textbf{worth testing}, \textbf{unclear} and
\textbf{unlikely}. The scoring is implemented as pure functions so that the web application and
the offline validation share one code path, which is what makes the validation below meaningful.

\subsection{Validation}

Per-protein ROC-AUC on scaffold-grouped held-out molecules (Table~\ref{tab:targets}) measures the
individual models. It does not measure the composition in Equation~\ref{eq:illness-score}, so two
further checks exercise the whole chain.

\paragraph{Repurposing retrieval.} Every illness in the knowledge base carries the small-molecule
drugs already approved for it. Those drugs are pushed through the complete pipeline --- featurise,
run every protein model, weight by association score --- and we record where their own indication
lands in the resulting ranking of all \TgtNDiseases{} illnesses. If the composition carries
signal, a drug's own illness should rank near the top; random ranking would place it first
\TgtValRandom\% of the time.

Over \TgtValDrugs{} approved drugs the pipeline puts the drug's own illness first \TgtValTopOne\%
of the time against a random \TgtValRandom\%, in its top three \TgtValTopThree\% of the time
against a random $20\%$, and at a median rank of \TgtValMedian{} of \TgtNDiseases{} where random
ranking would give $8.0$. Better than chance, then, but not by the margin the per-protein
ROC-AUCs might suggest --- and the aggregate is the least informative way to read this result.

\maybeinput{targets_validation}

\paragraph{The aggregate hides the actual finding.} The per-illness rows split into two groups
with nothing in between, and the dividing line is not model quality.

For the \TgtValApplicableIllnesses{} illnesses whose drugs act on \emph{human} proteins, retrieval
works: \TgtValApplicableTopThree\% top-three over \TgtValApplicableDrugs{} drugs at a median rank
of \TgtValApplicableMedian{}, roughly twice the random baseline, and far better than that at the
top --- every antidepressant but two puts depression in its own top three, and so does almost
every Alzheimer's drug.

For the \TgtValInfectiousIllnesses{} infectious diseases --- HIV infection, tuberculosis, malaria
and COVID-19 --- retrieval is \textbf{\TgtValInfectiousTopThree\% across all
\TgtValInfectiousDrugs{} drugs}, with a median rank of \TgtValInfectiousMedian{}, which is
\emph{worse} than random. Not one drug recovers its own indication.

That is not a modelling failure, and it is worth being precise about why. Open Targets associates
an illness with the \textbf{human} proteins implicated in it. But an antiretroviral inhibits HIV
protease or reverse transcriptase; an antimalarial acts on the parasite; a tuberculosis drug acts
on mycobacterial enzymes. None of those proteins is human, so none of them can appear in the
panel, and the models are being asked whether these drugs bind human proteins they were never
designed to bind. They correctly answer no. The pipeline is right about the chemistry and wrong
about the illness, because the knowledge base cannot express where those drugs act.

This bounds the module honestly: the method applies to illnesses treated by modulating a human
protein, and does not apply to infectious disease at all. A caveat discovered by measurement is
worth more than one asserted in advance, and it is the reason this validation was run on every
illness rather than on a favourable subset.

\paragraph{A deliberately permissive test.} The same drugs appear in ChEMBL, so a model may partly
be recalling them, and the scaffold-grouped per-protein figures in Table~\ref{tab:targets} remain
the strict measurement. Retrieval is reported because it is the only check that exercises the
composition end to end --- a pipeline that \emph{fails} it is broken regardless of its per-protein
scores --- and because, as above, the shape of the failures is more informative than the headline.

\paragraph{Negative controls.} Molecules nobody would mistake for drugs --- water, ethanol,
methane, glucose, acetic acid and benzene --- are run through the same chain. If they light up
every illness, the models are reading drug-likeness rather than protein activity. \textbf{Not one
of them is ever reported as \emph{worth testing}} --- \TgtCtrlFlags{} such verdicts across all six
control molecules and all \TgtNDiseases{} illnesses.

They are not dismissed outright either, and that is a real limitation of
Equation~\ref{eq:illness-score}: because the score is a maximum, it cannot fall below the single
most optimistic protein model, and these tiny molecules collect scores up to \TgtCtrlMax\% --- the
\emph{unclear} band rather than \emph{unlikely}. What keeps them out of the promising band is the
applicability-domain check: benzene resembles nothing in any training set of actives, confidence
comes back low, and the verdict is capped. The maximum and the domain check are therefore not
independent design choices --- the first is only safe because the second is there.

\subsection{Semi-supervised learning on a third domain}

Proteins with few measured molecules are precisely the case semi-supervised learning is supposed
to address, which makes this module a third independent test of the project's central protocol.
The five methods compared in Section~\ref{sec:methods} are rerun on the \TgtSSLNTargets{} proteins
with the fewest measured molecules, using the unlabelled ZINC pool --- the same experimental
design used for materials with GNoME and for toxicity with ZINC. The label budget is capped at
\TgtSSLBudget{}, but these proteins have so little measured data that only
\TgtSSLLabelsMin{}--\TgtSSLLabelsMax{} labels are actually available, which is the regime of
interest. Presumed inactives are excluded from the labelled pool here: they are assumed rather
than measured, and mixing them in would make the task artificially easy.

\maybeinput{targets_ssl}

The gains run from \TgtSSLMin\% to \TgtSSLMax\%, a mean of \TgtSSLMean\%, and on
\TgtSSLNoGain{} of the \TgtSSLNTargets{} proteins \textbf{no semi-supervised variant beat plain
supervised learning at all} --- the best method was the baseline. Where something did help it was
Self-Training on VDR and Co-Training on DPP9, the same two methods that won on materials and
toxicity.

This is a modest result, and it is the same modest result the other two domains gave: measured
across materials, toxicity and now protein activity, the honest summary of this project's central
experiment is that semi-supervised learning on an unlabelled pool buys a fraction of a percent
when it works and nothing when it does not. Reporting that across three independent domains is
worth more than reporting a larger number from one. It also agrees with the clearest negative
result in the materials track, where the GNoME pool actively \emph{hurt} magnetism prediction: a
large pool of unlabelled molecules is only useful when it is drawn from the same chemistry as the
labelled set, and ZINC is not where measured kinase or receptor ligands come from.

\subsection{What this module cannot do}

These limits are stated on every page of the application as well as here, and they are not
formalities --- each one marks a real boundary of the method.

\begin{itemize}
  \item A protein-level prediction says nothing about absorption, dose, selectivity, metabolism or
    safety in a person. Those are the reasons most candidate molecules fail, and none of them is
    modelled.
  \item Binding a protein associated with an illness is a \emph{hypothesis} about where to look
    next, not an effect. An association score records that researchers have linked a protein to an
    illness, not that acting on it helps.
  \item ChEMBL is biased toward published actives, and the presumed inactives used to compensate
    are an assumption. This is why the decoy fraction and the measured-only score are reported per
    model rather than averaged away.
  \item Many illnesses have no well-characterised druggable protein, so a low score may mean "no
    useful protein model exists for this illness" rather than "this molecule does nothing". That
    is exactly what the coverage figure is for.
\end{itemize}
